% ============================================================
% Plantilla: Artículo con normas Vancouver (ciencias de la salud)
% Las citas se numeran en orden de aparición: \cite{clave} -> [1]
% Recompilar ejecuta biber automáticamente.
% ============================================================
\documentclass[11pt,a4paper]{article}

\usepackage[utf8]{inputenc}
\usepackage[T1]{fontenc}
\usepackage{lmodern}
\usepackage[spanish,es-tabla]{babel}
\usepackage[margin=2.5cm]{geometry}
\usepackage{csquotes}
\usepackage{booktabs}
\usepackage[style=vancouver,backend=biber]{biblatex}
\usepackage[hidelinks]{hyperref}

\addbibresource{referencias.bib}

\title{Título del artículo de investigación en salud}
\author{Primer Autor\textsuperscript{1}, Segundo Autor\textsuperscript{2}\\[0.5em]
  \small\textsuperscript{1}Facultad de Ciencias Médicas, Universidad \quad
  \textsuperscript{2}Hospital de Referencia}
\date{\today}

\begin{document}

\maketitle

\begin{abstract}
\textbf{Objetivo:} indica el propósito del estudio.
\textbf{Métodos:} describe el diseño, la población y el análisis.
\textbf{Resultados:} presenta los hallazgos principales.
\textbf{Conclusiones:} resume la implicación clínica.

\medskip
\noindent\textbf{Palabras clave:} término MeSH 1; término MeSH 2; término MeSH 3.
\end{abstract}

\section{Introducción}
Presenta el problema de salud y los antecedentes \cite{icmje2019}. Las citas se
numeran en el orden en que aparecen en el texto \cite{lamport1994,knuth1984}.

\section{Métodos}
Describe el diseño del estudio, los criterios de inclusión y el análisis estadístico.

\section{Resultados}
\begin{table}[htbp]
  \centering
  \caption{Características de los participantes (n = 120).}
  \begin{tabular}{lc}
    \toprule
    Característica & Valor \\
    \midrule
    Edad, media (DE) & 45.2 (12.1) \\
    Mujeres, n (\%) & 64 (53.3) \\
    \bottomrule
  \end{tabular}
\end{table}

\section{Discusión}
Compara los resultados con estudios previos \cite{goossens2004}.

\printbibliography[title=Referencias]

\end{document}
